% !PRINTED 553
\weberpage{553}
%<<CONT
gemeinsames Element enthalten, vollständig identisch sein müssen, und die Classeneintheilung ist also durchaus eindeutig. Jede Classe ist durch ein beliebiges in ihr enthaltenes Functional, einen \gesperrt{Repräsentanten}, völlig bestimmt.

3. \gesperrt{Zwei mit einander associirte Functionale sind auch äquivalent und kommen daher in derselben Classe vor.}

Denn wenn $\varphi$ und $\psi$ associirt sind, so ist ihr Quotient $\varphi:\psi$ eine Einheit, und $\varphi$ und $\psi$ sind also auch äquivalent.

Eine Classe enthält aber nicht bloss die einzelnen Functionale $\varphi$, sondern alle durch diese Functionale bestimmten \gesperrt{Ideale}, und wir nennen diese Classen daher, wenn eine genauere Bezeichnung nöthig ist, \gesperrt{Functionalclassen} oder, häufiger noch, dem üblichen Sprachgebrauche gemäss, \gesperrt{Idealclassen}.

4. \gesperrt{Die Gesammtheit aller ganzen und gebrochenen Zahlen des Körpers $\Omega$, verbunden mit den sämmtlichen Einheiten und den Producten von Zahlen mit Einheiten, bilden unter sich eine Classe, die die Hauptclasse genannt wird.}

Als Repräsentanten der Hauptclasse kann man z. B. die Zahl 1 betrachten. Die Hauptclasse, als Idealclasse aufgefasst, enthält das Ideal $\mathfrak{o}$ und wird daher in der Folge durch den Buchstaben $O$ bezeichnet.

5. \gesperrt{In jeder Idealclasse giebt es ganze Functionale.}

Denn nach der Definition ist, wenn $\varphi$ irgend ein Functional und $\alpha$ eine Zahl ist, $\varphi$ mit $\alpha\varphi$ äquivalent. Wir können aber nach \S.\,137, 5. die Zahl $\alpha$, sogar rational, so bestimmen, dass $\alpha\varphi$ ein ganzes Functional wird.

6. \gesperrt{Aus jeder Idealclasse $C$ können wir einen Repräsentanten $\varphi$ auswählen, der nicht nur selbst ein ganzes Functional ist, sondern auch zu einem beliebig gegebenen ganzen Functional $\omega$ relativ prim ist.}

Nehmen wir, um diesen Satz zu beweisen, zunächst nach 5. einen beliebigen ganzen Repräsentanten $\varphi$ der Classe $C$ und eine durch $\varphi$ theilbare ganze Zahl $\alpha$, so ist

\weq{2}{\varphi\chi = \alpha,}
%CONT>>
